\section{Effective construction of the fixed rational spectral
target}\label{effective-construction-of-the-fixed-rational-spectral-target}

CC-XC11 identifies the rational isomorphism class of the target for a
smooth projective variety from its Frobenius polynomials. Here the
remaining choice of a concrete abelian-sector presentation is also made
effective. This construction always terminates on the stated spectral
input. The subsequent search for cross-cycles from the arbitrary variety
retains its one-sided termination condition.

The inputs from the established abelian sector are precisely Tate
surjectivity, numerical--homological equality, and Frobenius
semisimplicity on abelian varieties and their products, as recorded in
AQ-ABELIAN-RECOGNITION and CC-TATE01--02. The algorithm never assumes
any of these properties for the arbitrary input variety.

\subsection{CC-XC14: a terminating concrete-target
algorithm}\label{cc-xc14-a-terminating-concrete-target-algorithm}

Fix \(q=p^a\). Suppose a finite list \[
 P_0(t),\ldots,P_N(t)\in\mathbf Z[t]                       \tag{CC-XC.E1}
\] has the following properties. Each polynomial is monic, has nonzero
constant term, and all its roots are integral \(q\)-Weil numbers of
weight equal to its index. The polynomial \(1\), with no roots, is
permitted. For every rational irreducible factor \(f^{m_f}\) of \(P_i\),
the Schur index \(d_f\) from (CC-XC.R15) divides \(m_f\). These are
promised conditions on the input; they are satisfied by the Frobenius
polynomials of every smooth projective \(X/\mathbf F_q\), by CC-XC11,
and also by the polynomials of a proper smooth scheme supplied as input,
by CC-XC13.

If a smooth scheme has components of different dimensions, split it into
the finitely many open-and-closed unions of components of each fixed
dimension. These unions are defined over \(k\), and their degreewise
polynomials multiply to the polynomials of the whole scheme. The index
\(d_f\) depends only on \(f,q,i\), so its divisibility survives addition
of the factor multiplicities. The final cross-cycle procedure below is
applied separately to these pure-dimensional pieces; the whole object
belongs to the abelian sector exactly when each piece does, since that
sector is closed under finite direct sums and direct summands.

\textbf{Theorem.} There is a terminating algorithm that constructs
finitely many explicitly presented abelian varieties and finite étale
schemes \(Y_\nu/\mathbf F_q\), together with a rational matrix of
algebraic cycles representing a numerical idempotent \[
 e\in\operatorname{End}_{\mathcal A_q}(T),\qquad
 T=\bigoplus_\nu h(Y_\nu),                               \tag{CC-XC.E2}
\] such that, for every \(\ell\ne p\), \[
 \det(t-F\mid R_\ell(T,e)^i)=
 \begin{cases}
 P_i(t),&0\le i\le N,\\
 1,&i>N.
 \end{cases}                                           \tag{CC-XC.E3}
\] Here \(R_\ell(T,e)^i\) is the degree-\(i\) realization; the projector
is valid in every such realization. The object \((T,e)\) is the unique
rational spectral target specified by the list.

For a smooth projective input variety \(X/\mathbf F_q\) given by
equations, one may first compute its Frobenius polynomials and then run
this algorithm. After these two terminating steps, the fixed-target
cross-cycle search in CC-XC12 halts exactly when the cohomological
motive of \(X\) belongs to the abelian sector. This conclusion does not
assert termination of that last search for every \(X\), and gives no
complexity or cycle-degree bound.

\subsubsection{Complete proof and
algorithm}\label{complete-proof-and-algorithm}

\textbf{1. Computational operations and numerical ranks on abelian
products.} A smooth projective geometrically integral variety over a
finite field is called a ``nice variety'' in Poonen--Testa--van Luijk
(2015). Their Proposition 8.1, author p.~10, computes its zeta function.
Separating its factors by their weights then computes the degreewise
Frobenius polynomials, as in the proof of Proposition 8.2 on that page.
Every abelian variety, including every product used below, meets these
hypotheses.

Let \(Z/k\) be an abelian variety of dimension \(D\), and fix
\(0\le r\le D\). Write \(\operatorname{Num}^r(Z)_{\mathbf Q}\) for
rational codimension-\(r\) cycles over \(k\) modulo numerical
equivalence. Its dimension is the computable integer \[
 R_r(Z)=
 \operatorname{ord}_{t=q^r}
       \det(t-F\mid H^{2r}(Z_{\bar k},\mathbf Q_\ell)).
                                                        \tag{CC-XC.E4}
\] Indeed, semisimplicity identifies the right side with the dimension
of the Frobenius-invariant subspace of \(H^{2r}(Z)(r)\). Tate
surjectivity and numerical--homological equality, with the
scalar-injectivity statement in AQ.15, identify that space with
\(\operatorname{Num}^r(Z)_{\mathbf Q}\otimes\mathbf Q_\ell\). Only the
known abelian sector is used in this rank computation.

The complementary intersection pairing between
\(\operatorname{Num}^r(Z)_{\mathbf Q}\) and
\(\operatorname{Num}^{D-r}(Z)_{\mathbf Q}\) is perfect. By the
definition of numerical equivalence it is nondegenerate on both sides;
the spaces are finite-dimensional, so it is a perfect pairing and their
dimensions agree. Enumerate \(R=R_r(Z)\) cycles \(z_1,\ldots,z_R\) of
codimension \(r\) and \(R\) cycles \(w_1,\ldots,w_R\) of complementary
codimension until \[
             \det(z_u\cdot w_v)_{1\le u,v\le R}\ne0.
                                                        \tag{CC-XC.E5}
\] Such a pair of lists exists because cycles over \(k\) span both
numerical spaces. Each candidate is finitely checkable:
Poonen--Testa--van Luijk, Lemma 8.7, author p.~11, computes the
intersection numbers. Exhaustive enumeration therefore terminates. When
\(R=0\), use empty lists.

The classes of the \(z_u\) form a rational basis. For any further
rational cycle \(z\), the intersection numbers \(z\cdot w_v\), followed
by inversion of the rational matrix in (CC-XC.E5), give its coordinates
in that basis. In particular, numerical equality is now decidable for
cycles of this codimension on \(Z\). The argument is the
finite-dimensional intersection-matrix method appearing in
Poonen--Testa--van Luijk, Theorem 8.15(b), author p.~13, specialized
here to a known arithmetic numerical rank. It does not call their
conditional algorithm on an arbitrary variety.

\textbf{2. Effective numerical endomorphisms of an abelian sum.} Let \[
                   T_+=\bigoplus_{\nu=1}^{s}h(A_\nu)
                                                        \tag{CC-XC.E6}
\] be a finite sum of abelian varieties over \(k\). For every ordered
pair \((\nu,\mu)\), apply Step 1 to \[
 Z=A_\nu\times A_\mu,\qquad r=\dim A_\nu.
 \] This computes a rational cycle basis for the block
\(\operatorname{Hom}(h(A_\nu),h(A_\mu))\) in the numerical category,
with a finite coordinate test for every cycle in the block. Taking all
blocks gives a rational basis for the finite-dimensional algebra
\(\operatorname{End}(T_+)\).

Composition is also computable. Represent the two correspondences by
cycles, pull them to the triple product, intersect, and push forward to
the first and third factors. All these varieties are smooth projective.
The effective Chow-intersection operation is Poonen--Testa--van Luijk,
Remark 8.18, author p.~14; the pullbacks along projections and proper
pushforwards are computable from equations and generic degrees, as in
the complete verifier CC-XC05. Step 1 gives the numerical coordinates of
the resulting cycle. Thus for each rational element \(e\) of this
algebra the test \[
                              e^2=e                     \tag{CC-XC.E7}
\] is a finite exact rational calculation. No test of numerical
equivalence on a product involving the arbitrary \(X\) is being used.

\textbf{3. Testing the image polynomial of a projector.} Suppose \(e\)
passes (CC-XC.E7). Numerical--homological equality on the abelian
products makes its realization an idempotent in every \(\ell\)-adic
realization. All cycles are defined over \(k\); consequently this
idempotent commutes with Frobenius.

Compute the degreewise Frobenius polynomials of every \(A_\nu\), and the
rational cohomological projectors \(\pi_{i,\nu}\) by the Chinese
remainder construction in CC-XC05. Here \(F_\nu\) is the correspondence
acting as geometric Frobenius, represented by the transpose graph of the
relative \(q\)-power endomorphism in the cohomological correspondence
convention. For \(n\ge0\), the rational number \[
 s_{i,n}=(-1)^i\sum_\nu
   \operatorname{Tr}_{\mathrm{corr}}
       \bigl(\pi_{i,\nu}F_\nu^n e_{\nu\nu}\bigr)
                                                        \tag{CC-XC.E8}
\] is computable by intersection with the diagonal. The notation
\(\operatorname{Tr}_{\mathrm{corr}}\) is the ordinary intersection
trace, which is the alternating cohomological trace. Hence the factor
\((-1)^i\) in (CC-XC.E8) gives exactly \[
 s_{i,n}=\operatorname{Tr}
       (F^n\mid\operatorname{im}e\cap H^i(T_+)).
 \] For the direct sum the trace is the sum of the diagonal block
traces; off-diagonal blocks do not contribute. The formula holds at
every \(\ell\ne p\) and is independent of it.

In particular \(s_{i,0}\) is the integer rank of the image in degree
\(i\). First test that it equals the desired degree of \(P_i\), using
the value zero outside the specified degrees. For a matched rank
\(b_i\), the \(b_i\) further traces \(s_{i,1},\ldots,s_{i,b_i}\)
determine the characteristic polynomial of Frobenius on that image by
Newton's identities. Compare it with \(P_i\). Test every degree up to
the maximum of \(N\) and twice the largest dimension of the \(A_\nu\).
The realization is zero above the latter bound, so any nonconstant
prescribed polynomial in that range rejects the candidate. All larger
degrees vanish on both sides. This gives a finite exact test of every
required image polynomial. For a zero desired rank, the rank test
suffices.

\textbf{4. Explicit degree-zero target.} Every integral weight-zero Weil
number is a root of unity, by Kronecker's theorem, so factor \[
                       P_0(t)=\prod_n\Phi_n(t)^{b_n},
                                                        \tag{CC-XC.E9}
\] with finitely many nonzero \(b_n\). One can find the indices \(n\) by
enumerating cyclotomic polynomials and matching the irreducible factors;
the root-of-unity property guarantees termination. For each \(n\), the
Frobenius on \(h(\operatorname{Spec}\mathbf F_{q^n})\) is a regular
cyclic permutation with characteristic polynomial \(t^n-1\). Choose by
the rational Chinese remainder algorithm a polynomial \(u_n(t)\) that is
\(1\) modulo \(\Phi_n(t)\) and \(0\) modulo every other cyclotomic
factor of \(t^n-1\).

The explicit rational correspondence \(u_n(F)\), a rational linear
combination of Frobenius graphs, is an idempotent on this finite étale
scheme and its image has polynomial \(\Phi_n(t)\). The identity holds as
a correspondence here: \(F^n=1\) and the permutation action of
\(\mathbf Q[t]/(t^n-1)\) is faithful. Take \(b_n\) copies. This
constructs a concrete rational Artin projector \((T_0,e_0)\) with
polynomial \(P_0\). It is zero if \(P_0=1\).

\textbf{5. Exhaustive search and termination.} Use the effective
enumeration of abelian varieties and finite lists established in
CC-XC07: enumerate projective equations together with explicit identity,
inverse and group law on finite affine covers, and retain the smooth
geometrically connected commutative group schemes. The finite algebraic
verifications and their exhaustiveness are explained in that proof. For
each finite list \(T_+\), Steps 1--2 terminate and give a finite
rational basis of \(\operatorname{End}(T_+)\). Enumerate all rational
coefficient tuples in that basis, retain the idempotents, and apply Step
3. The desired positive-degree list is \(P_i\) for \(1\le i\le N\), with
polynomial \(1\) in degree zero and outside that range. Dovetail these
searches over all lists so that each finite stage in each search is
eventually executed.

The promised index divisibilities give the unique rational object
\(\bigoplus_{i\ge1,f}S_f^{m_f/d_f}\) by CC-AG08. The
effective-presentation argument in CC-XC11 proves that it is a rational
summand of a finite sum of abelian varieties: take a weight-one root for
each positive weight eigenvalue, apply Honda, and take tensor powers and
enough copies. This argument needs only the promised spectral
conditions, not the existence of an \(X\). Consequently some enumerated
list contains it, and some rational idempotent in the computed basis has
exactly the required image polynomials. The dovetailed search eventually
accepts this projector.

Return \[
                     (T,e)=(T_0\oplus T_+,e_0\oplus e).
                                                        \tag{CC-XC.E10}
\] The construction satisfies (CC-XC.E3), and uniqueness follows from
CC-AG08. If every positive-degree polynomial is \(1\), this step can
instead return \(T_+=0\) immediately. The completely zero list returns
the zero object.

Finally, if a smooth projective variety \(X\) is the original input, its
Frobenius polynomials are computable by CC-XC05 using Poonen--Testa--van
Luijk, Propositions 8.1--8.2. CC-XC11 verifies all the promises, so the
construction just proved terminates and supplies the concrete target
required by CC-XC12. That theorem then proves the exact one-sided
termination assertion for the remaining cross-cycle search. \(\square\)

\subsection{Scope of the algorithmic
conclusion}\label{scope-of-the-algorithmic-conclusion}

The guaranteed stage uses cycle searches only on products of abelian
varieties, where the relevant ranks and equivalence relations are
already known. It produces equations and rational projector cycles,
rather than only a list of Weil labels. The search on products involving
\(X\) begins only after that target has been fixed, and its universal
termination is not proved. In particular this theorem does not assert an
algorithm deciding arbitrary Tate classes or numerical equivalence on
arbitrary smooth projective varieties.

When proper nonprojective smooth cohomology polynomials are supplied as
spectral input, the target construction still applies by CC-XC13. The
final assertion about computing from projective equations and applying
the fixed-target cycle verifier retains its projective scope.

\subsection{APA bibliography and exact
locators}\label{apa-bibliography-and-exact-locators}

Poonen, B., Testa, D., \& van Luijk, R. (2015). Computing Néron--Severi
groups and cycle class groups. \emph{Compositio Mathematica, 151}(4),
713--734. https://doi.org/10.1112/S0010437X14007878 Inspected author
manuscript:
https://math.mit.edu/\textasciitilde poonen/papers/compute\_ns.pdf Exact
author-page locators: §7.1, pp.~6--7; Propositions 8.1--8.2, p.~10;
Lemma 8.7, p.~11; Theorem 8.15(b), p.~13; Remark 8.18, p.~14. Theorem
8.36(b), p.~19, has its own conditional hypotheses and is not used on
arbitrary \(X\) in this construction.

Honda, T. (1968). Isogeny classes of abelian varieties over finite
fields. \emph{Journal of the Mathematical Society of Japan, 20}(1--2),
83--95. Main Theorem, p.~83. https://doi.org/10.2969/jmsj/02010083

Milne, J. S. (2020). \emph{Algebraic number theory} (Version 3.08, 19
July 2020). Corollary 5.6, printed p.~87, for the Kronecker
root-of-unity criterion. https://www.jmilne.org/math/CourseNotes/ANT.pdf

The exact all-abelian primary inputs and their scope are those in
CC-TATE01--02 and AQ.15. The complete rational classification, cycle
enumeration, and intersection-trace proofs are CC-AG08, CC-XC05--07, and
CC-XC11--13.
